% SOURCES: README.md; docs/1.0.0v/supervisor-report.md (§1, §13); draft/Features/features.md
% Abstract — written last, from the verified results in ch. 11. The pending
% render-layer performance figures (issue #95) are deliberately NOT claimed here.

Authoring a web page that is at once \emph{semantic}, \emph{responsive}, and
\emph{accessible} still demands specialist skill, and the visual ``no-code''
builders that promise to remove that barrier typically do so at the expense of
output quality: heavy proprietary runtimes, non-semantic markup, absolute
positioning, and pages that are difficult to audit for accessibility.
\textbf{Draw-to-Web} is a desktop application that resolves this tension. A single
\emph{Document Model} is the only source of truth: the canvas is a rendering of
that model, and a deterministic generator walks the same model to emit semantic
HTML5, tokens-driven CSS, and opt-in, independently toggleable JavaScript. When no
runtime behaviour is enabled, the output contains no \texttt{<script>} at all.
Accessibility is enforced as a build-time hard gate --- \texttt{axe-core} runs on
the generated HTML and blocks export on any \emph{critical} or \emph{serious}
violation --- rather than left to manual review. The system was built as a
strict-TypeScript Electron + React application with a Zustand/immer state layer, a
Zod-validated process boundary, and a nine-stage export pipeline. Version~1.0.0
adds three authoring surfaces that reuse the same operations and export path:
\emph{draw-to-create} (draw a rectangle to place an element), \emph{match-layout}
(adopt a professional design from a rough draft), and a headless \emph{Model
Context Protocol} (MCP) server. The v1.0.0 release is green across all automated
gates: \textbf{1042 automated tests} across 110 files pass, linting and strict
type-checking are clean, the accessibility gate is enforced on every export path,
and the data-layer and export performance budgets are met with wide margins (for
example a full nine-stage export of a portfolio page completes in roughly
0.2--0.3~s against a 10~s budget). The work shows that visual, no-code authoring
and hand-written-quality output need not be opposites.

\vspace{0.6em}
\noindent\textbf{Keywords:} visual web authoring; no-code; semantic HTML5;
web accessibility (WCAG); design tokens; deterministic code generation; Electron;
Model Context Protocol.
